\chapter{Configuration}
\label{chap:config}

La configuration (\fichier{config.py}) est pensée pour être écrite par un aidant qui n'est pas programmeur. D'où trois principes :
\begin{itemize}
  \item tous les réglages sont optionnels ;
  \item les erreurs sont détectées \emph{au chargement}, toutes à la fois, plutôt qu'au premier clic ;
  \item chaque message en français désigne l'endroit exact du problème et propose une correction.
\end{itemize}

\section{Formats}

Le format se déduit de l'extension : \fichier{.yaml}/\fichier{.yml} (PyYAML, \code{safe_load}, qui ne construit pas d'objets Python arbitraires), \fichier{.toml} (\code{tomllib}, bibliothèque standard depuis Python 3.11) ou \fichier{.json}. La structure est la même dans les trois formats ; \fichier{victor/exemples/} contient la configuration par défaut en TOML et en JSON.

\begin{itemize}
  \item Le fichier est lu en UTF-8 avec \code{utf-8-sig}, qui ignore la marque d'ordre des octets (BOM) qu'ajoutent certains éditeurs sous Windows.
  \item Une erreur de syntaxe donne la ligne et la colonne (YAML, JSON), ou le message du lecteur TOML, ainsi qu'un conseil propre au format : indentation par espaces en YAML, guillemets en TOML, virgule finale interdite en JSON.
  \item Si le fichier n'existe pas, toutes les valeurs par défaut s'appliquent et un message le signale.
\end{itemize}

\section{Fusion avec les valeurs par défaut}

\code{_fusionner(defaut, donne, chemin, erreurs)} parcourt récursivement le dictionnaire \code{DEFAUTS}. \emph{Le type de la valeur par défaut détermine la validation} :

\begin{table}[H]
\centering\small
\begin{tabularx}{\textwidth}{@{}l L@{}}
\toprule
Type par défaut & Règle \\
\midrule
dictionnaire & fusion récursive ; une valeur non dictionnaire est une erreur \\
booléen & \code{oui}/\code{vrai}/\code{o}/\code{true} ou \code{non}/\code{faux}/\code{n}/\code{false}, sans distinction de casse \\
nombre & doit être un nombre (un booléen n'est pas accepté), dans les bornes \code{BORNES} si elles existent \\
liste & même longueur que la valeur par défaut, uniquement des nombres \\
texte & converti en chaîne ; s'il figure dans \code{CHOIX}, doit être une valeur autorisée \\
\bottomrule
\end{tabularx}
\caption{Règles de validation de \code{_fusionner}.}
\end{table}

Une clé inconnue produit une erreur avec la clé connue la plus proche (\code{difflib.get_close_matches}, fondé sur la similarité de Ratcliff-Obershelp), par exemple \enquote{\code{reglages.duree_clik} : réglage inconnu, voulais-tu dire \code{"duree_clic"} ?}. Une valeur invalide est remplacée par sa valeur par défaut \emph{et} signalée : la vérification continue, ce qui permet d'afficher toutes les erreurs en un seul passage.

\paragraph{Exclusion des booléens.} En Python, \code{True} est un \code{int}. Sans le test \code{isinstance(v, bool)}, \code{duree_clic: oui} serait accepté comme la valeur 1.

\section{Conversions spécifiques}

\begin{itemize}
  \item \code{grille} : l'expression régulière \verb|(\d+)\s*[xX]\s*(\d+)| (ainsi que le signe $\times$) accepte \enquote{2x3}, \enquote{2 X 3} ou \enquote{2×3}. Lignes et colonnes doivent être entre 1 et 4 : au-delà, les cases deviennent trop petites pour la précision d'un eye tracker webcam (\empi).
  \item \code{couleur_accent} et couleurs de boutons : un nom parmi 10 (rouge, vert, bleu, orange, jaune, violet, rose, gris, turquoise, marron) ou \texttt{\#RRGGBB}. Les valeurs hexadécimales des noms correspondent à une palette d'interface aux teintes vives, lisibles sur fond sombre comme sur fond clair (\conv).
\end{itemize}

\section{Menus et boutons}

\code{_lire_menus} construit un dictionnaire \code{nom -> liste de Bouton} :
\begin{enumerate}
  \item Le menu \code{principal} est obligatoire dès qu'il y a des menus ; s'il n'y a aucun menu, un menu principal vide est créé.
  \item Chaque sous-menu reçoit \emph{en première position} un bouton Retour (\code{texte_retour}, action \code{retour}, \code{est_retour=True}). Il apparaît donc toujours dans la première case de la première page, à la même place dans tous les sous-menus.
  \item Clés d'un bouton : \code{texte}, \code{action}, \code{icone}, \code{duree_clic}, \code{couleur}. Toute autre clé est une erreur, avec suggestion.
  \item \code{action} est une chaîne ou une liste de chaînes. Chacune est analysée par \code{actions.analyser}, qui reçoit la liste de \emph{tous} les noms de menus : un \code{ouvrir} vers un menu défini plus bas dans le fichier est accepté, une faute de frappe est détectée.
  \item \code{icone} est un chemin relatif au fichier de configuration ; l'existence du fichier est vérifiée.
  \item \code{duree_clic} doit être entre \qtyrange{0.2}{30}{\second}.
  \item Il faut au moins un texte ou une icône.
\end{enumerate}
Les messages localisent l'erreur, par exemple \enquote{\texttt{menus.Raccourcis, bouton 2 «~Navigateur~» : ...}}.

\paragraph{Le type \code{Bouton}.} C'est une \code{dataclass} avec \code{texte}, \code{etapes} (liste de \code{(nom, [arguments])}), \code{icone} (chemin absolu), \code{duree_clic}, \code{couleur} (RGB), \code{est_retour} et \code{navigue}. Ce dernier champ est calculé : vrai si la séquence contient \code{ouvrir}, \code{retour} ou \code{accueil}. Il sert à ne pas afficher de bulle pour une simple navigation.

\section{Résultat}

S'il y a au moins une erreur, \code{charger} lève une unique \code{ErreurConfig} qui les liste toutes. Sinon, elle renvoie un \code{SimpleNamespace} imbriqué (\code{cfg.ecran.zone_morte.largeur}, \code{cfg.avance.camera}…), dont l'attribut \code{menus} est le dictionnaire des boutons.

\section{Tous les réglages}

\begin{longtable}{@{}p{5.1cm} p{1.9cm} p{1.6cm} p{1.9cm} p{3.8cm}@{}}
\toprule
Clé & \fichier{config.py} & \fichier{config.yaml} & Bornes & Effet \\
\midrule
\endfirsthead
\toprule
Clé & \fichier{config.py} & \fichier{config.yaml} & Bornes & Effet \\
\midrule
\endhead
\bottomrule
\endfoot
\code{ecran.plein_ecran} & non & oui & & plein écran au démarrage \\
\code{ecran.taille_fenetre} & [1280, 720] & & 2 nombres & taille hors plein écran \\
\code{ecran.grille} & 2x2 & & 1--4 & lignes × colonnes \\
\code{ecran.largeur_bandes} & 7 & 12 & 0--30 \% & largeur de chaque bande \\
\code{ecran.zone_morte.largeur} & 22 & 30 & 0--80 \% & largeur totale \\
\code{ecran.zone_morte.hauteur} & 34 & 40 & 0--80 \% & hauteur totale \\
\code{ecran.zone_morte.aplatissement_haut} & \num{0.5} & 0 & 0--1 & forme du haut \\
\code{ecran.theme} & sombre & & sombre, clair & thème \\
\code{ecran.couleur_accent} & \#4F8EF7 & & & accent \\
\code{ecran.taille_texte} & 100 & & 30--300 \% & taille du texte \\
\code{ecran.police} & (système) & & & fichier \fichier{.ttf} \\
\code{ecran.animations} & oui & & & glissement de page \\
\code{reglages.duree_clic} & \qty{2.0}{\second} & & 0,2--30 & durée de fixation \\
\code{reglages.retour_apres_action} & non & & & retour à l'accueil \\
\code{reglages.texte_retour} & $\leftarrow$ Retour & & & texte du bouton retour \\
\code{reglages.duree_message} & \qty{1.5}{\second} & & 0--30 & durée des bulles \\
\code{avance.camera} & 0 & & & index de la webcam \\
\code{avance.resolution_camera} & [640, 480] & & 2 nombres & résolution demandée \\
\code{avance.afficher_camera} & non & & & vignette de debug \\
\code{avance.images_par_seconde} & 30 & & 5--120 & cadence de l'interface \\
\code{avance.seuil_clignement} & \num{0.6} & & 0--1 & clignements \\
\code{avance.filtre_stabilite} & \num{0.6} & & 0,01--10 & $f_{\min}$ One Euro (Hz) \\
\code{avance.filtre_reactivite} & \num{3.0} & & 0--100 & $\beta$ One Euro \\
\code{avance.delai_perte_regard} & \qty{0.5}{\second} & & 0--10 & regard perdu \\
\code{avance.recalibration_clic} & oui & & & recalibration au clic \\
\code{avance.recalibration_part_centre} & \num{0.5} & \num{0.65} & 0--1 & $\alpha$ \\
\code{avance.recalibration_poids} & \num{1.0} & & 0--20 & $W$ \\
\code{avance.recalibration_oubli} & \num{0.8} & \num{0.85} & 0--1 & $d$ \\
\code{avance.recalibration_points_max} & 20 & & 1--200 & fenêtre glissante \\
\code{avance.recalibration_fenetre} & \qty{0.5}{\second} & \qty{0.4}{\second} & 0,05--5 & $\Delta$ \\
\end{longtable}

\begin{attention}
L'en-tête de \fichier{config.yaml} affirme que \enquote{les valeurs écrites ici SONT les valeurs par défaut}. C'est faux pour huit réglages, ceux dont la colonne \fichier{config.yaml} est remplie ci-dessus. Supprimer l'une de ces lignes du YAML ne rétablit donc pas la valeur qui y était écrite. Il faudrait aligner \code{DEFAUTS} sur le YAML, ou corriger le commentaire. De même, les constantes du module \fichier{utils_v2.py} (\code{RECAL_DECAY = 0.8}, \code{RECAL_WINDOW_S = 0.5}) ne correspondent qu'aux défauts de \fichier{config.py}.
\end{attention}
